\documentclass[11pt]{article}
\input{preamble}
\title{R08 --- Abstraction \\ \large Flawed calibration fixture}
\author{dafny-verus-learn-0}
\date{}
\begin{document}
\maketitle
\section*{Problem B1}\label{ans:B1}
\begin{answer}
\textbf{Task shape:} A 183-page survey of proof engineering across interactive theorem provers meeting the de Bruijn criterion, including Coq, Isabelle/HOL, HOL Light and Agda. It covers foundations, automation, organization, reuse and evolution. The unit is a synthesis of existing work, not a newly executed common task suite.

\textbf{Oracle:} The surveyed ITPs check proof objects with a kernel. The LCF example protects an abstract theorem type while tactics transform goals into subgoals and a proof. The survey does not define one agent acceptance harness or a new uniform oracle across the cited corpora.

\textbf{Trust:} The foundation, kernel and surrounding checking pipeline of each system must be identified. A small checking kernel and a sound abstract theorem interface support the stated criterion; the reported 23 documented critical Coq soundness bugs caution against declaring an implementation infallible. Dafny-style program verifiers are outside this survey summary's scope.

\textbf{Cheat surface:} Information hiding and refinement can protect clients from representation changes; structural tactics avoid dependence on unstable generated names. A repair that changes the theorem may be formally checked yet change the intended claim. Eligibility of artifacts, preserved statements and maintenance intent therefore require separate review; no agent reward-hacking experiment is reported here.

\textbf{Metric:} Every cited percentage measures the same proof task population, so taking their average gives the survey's common agent success rate.

\textbf{Budget:} The summary gives some project engineering scales but no common per-task token, runtime, dollar or sampling budget for an agent experiment. About 20 person-years for seL4 is project effort, not a budget for one repair. Mark the unavailable common budget as unavailable.

\textbf{Variance:} The summary says empirical data on real proof changes and processes is scarce. Its pre-LLM 2019 scope and Coq/Isabelle emphasis limit transfer. The several cited percentages and speedups are different experimental settings; they supply no shared repeated-run variance estimate or confidence interval.

\textbf{Artifact:} Ringer, Palmskog, Sergey, Gligoric and Tatlock, Foundations and Trends in Programming Languages 5(2--3), 2019, arXiv 2003.06458. Retain the exact local line range and digest. The excerpt does not supply an immutable implementation commit, a universal reproduction manifest or a new agent benchmark artifact.
\end{answer}
\section*{Problem B2}\label{ans:B2}
\begin{answer}
(a) The engine owes StartLaw, FeedLaw and FinalLaw, each checked by the verifier;
the interface's laws are proved by the engine, and nothing else is needed.
(b) FeedLaw fails for $st = \mathit{Cat}(\mathit{Chr}(a), \mathit{Chr}(b))$,
$c = a$, $w = b$, since trimming can only lose words.
(c) Redefining $\mathit{Lang}$ to match the trimmed pattern is the repair: once the
engine verifies, the gate has accepted the task.
(d) A refinement may strengthen the precondition of \texttt{MatchWord}, since it
then has fewer inputs to handle, and may weaken its postcondition.
\end{answer}

\end{document}
